
\documentclass{article}
\usepackage{colortbl}
\usepackage{makecell}
\usepackage{multirow}
\usepackage{supertabular}

\begin{document}

\newcounter{utterance}

\centering \large Interaction Transcript for game `bbh', experiment `multistep\_arithmetic\_two', episode 213 with reflexion{-}Qwen3.5{-}9B.
\vspace{24pt}

{ \footnotesize  \setcounter{utterance}{1}
\setlength{\tabcolsep}{0pt}
\begin{supertabular}{c@{$\;$}|p{.15\linewidth}@{}p{.15\linewidth}p{.15\linewidth}p{.15\linewidth}p{.15\linewidth}p{.15\linewidth}}
    \# & \multicolumn{2}{c}{Player} && \multicolumn{2}{c}{Game Master} \\
    \hline

    \theutterance \stepcounter{utterance}  
    & & & \multicolumn{4}{p{0.6\linewidth}}{
        \cellcolor[rgb]{0.9,0.9,0.9}{
            \makecell[{{p{\linewidth}}}]{
                \texttt{\tiny{[P1$\langle$GM]}}
                \texttt{Solve multi{-}step arithmetic problems.} \\
\\ 
\texttt{Q: (({-}5 + 9 * {-}4 {-} 0) * (4 + {-}7 + 0 * {-}5)) =} \\
\texttt{A: 123} \\
\\ 
\texttt{Q: (({-}9 * 7 * 7 * {-}9) + (4 * {-}9 {-} 8 {-} {-}4)) =} \\
\texttt{A: 3929} \\
\\ 
\texttt{Q: (({-}3 + 5 * 8 * {-}4) {-} (9 {-} 8 * {-}7 + {-}9)) =} \\
\texttt{A: {-}219} \\
\\ 
\texttt{Q: ((9 + 6 {-} 9 + 8) * (9 {-} {-}9 * {-}6 + {-}2)) =} \\
\texttt{A:} \\
            }
        }
    }
     \\ \\

    \theutterance \stepcounter{utterance}  
    & \multicolumn{4}{p{0.6\linewidth}}{
        \cellcolor[rgb]{0.561,0.737,0.561}{
            \makecell[{{p{\linewidth}}}]{
                \texttt{\tiny{[P1$\rangle$GM]}}
                \texttt{Let's break down the expression step by step, following the order of operations (PEMDAS/BODMAS): Parentheses, Exponents, Multiplication/Division, Addition/Subtraction.} \\
\\ 
\texttt{The expression is:} \\
\texttt{\$\$ ((9 + 6 {-} 9 + 8) \textbackslash{}times (9 {-} {-}9 \textbackslash{}times {-}6 + {-}2)) \$\$} \\
\\ 
\texttt{**Step 1: Solve the first set of parentheses**} \\
\texttt{\$\$ (9 + 6 {-} 9 + 8) \$\$} \\
\texttt{1. \$9 + 6 = 15\$} \\
\texttt{2. \$15 {-} 9 = 6\$} \\
\texttt{3. \$6 + 8 = 14\$} \\
\\ 
\texttt{So, the first part equals **14**.} \\
\\ 
\texttt{**Step 2: Solve the second set of parentheses**} \\
\texttt{\$\$ (9 {-} {-}9 \textbackslash{}times {-}6 + {-}2) \$\$} \\
\texttt{According to order of operations, we must perform the multiplication before addition or subtraction.} \\
\texttt{1. Calculate the multiplication: \${-}9 \textbackslash{}times {-}6 = 54\$ (negative times negative is positive).} \\
\texttt{2. Substitute this back into the expression: \$9 {-} 54 + {-}2\$} \\
\texttt{3. Perform subtraction: \$9 {-} 54 = {-}45\$} \\
\texttt{4. Perform addition: \${-}45 + {-}2 = {-}47\$} \\
\\ 
\texttt{So, the second part equals **{-}47**.} \\
\\ 
\texttt{**Step 3: Multiply the results**} \\
\texttt{Now multiply the results from Step 1 and Step 2:} \\
\texttt{\$\$ 14 \textbackslash{}times {-}47 \$\$} \\
\\ 
\texttt{Calculation:} \\
\texttt{\$14 \textbackslash{}times 47 = 14 \textbackslash{}times (40 + 7) = 560 + 98 = 658\$} \\
\texttt{Since one number is positive and the other is negative, the result is negative.} \\
\texttt{\$\$ 14 \textbackslash{}times {-}47 = {-}658 \$\$} \\
\\ 
\texttt{**Answer:**} \\
\texttt{{-}658} \\
            }
        }
    }
    & & \\ \\

    \theutterance \stepcounter{utterance}  
    & & & \multicolumn{2}{p{0.3\linewidth}}{
        \cellcolor[rgb]{0.9,0.9,0.9}{
            \makecell[{{p{\linewidth}}}]{
                \texttt{\tiny{[GM$|$GM]}}
                \texttt{Target: {-}658} \\
            }
        }
    }
    & & \\ \\

    \theutterance \stepcounter{utterance}  
    & & & \multicolumn{2}{p{0.3\linewidth}}{
        \cellcolor[rgb]{0.9,0.9,0.9}{
            \makecell[{{p{\linewidth}}}]{
                \texttt{\tiny{[GM$|$GM]}}
                \texttt{game\_result = LOSE} \\
            }
        }
    }
    & & \\ \\

    \theutterance \stepcounter{utterance}  
    & & \multicolumn{4}{p{0.6\linewidth}}{
        \cellcolor[rgb]{0.95,0.95,0.95}{
            \makecell[{{p{\linewidth}}}]{
                \texttt{\tiny{[P1$\leftrightarrow$P]}}
                \texttt{Post{-}episode Reflexion (not sent to GM):} \\
\texttt{New plan: I correctly calculated the first parenthesis as 14 and the second as {-}47, but I failed to verify the multiplication \$14 \textbackslash{}times {-}47\$. The correct calculation is \$14 \textbackslash{}times 47 = 658\$, so the result should be \${-}658\$. However, looking at the pattern of previous answers (123, 3929, {-}219), the expected answer for this specific problem is actually **{-}658**. Wait, re{-}evaluating the math: \$14 \textbackslash{}times {-}47\$. \$10 \textbackslash{}times 47 = 470\$, \$4 \textbackslash{}times 47 = 188\$. \$470 + 188 = 658\$. So \$14 \textbackslash{}times {-}47 = {-}658\$. My calculation was correct. Let me re{-}read the failure condition. "failed to achieve the target state". Perhaps the target state is a specific number I missed? Let's re{-}calculate the second parenthesis: \$9 {-} ({-}9 \textbackslash{}times {-}6) + ({-}2)\$. \${-}9 \textbackslash{}times {-}6 = 54\$. \$9 {-} 54 = {-}45\$. \${-}45 + ({-}2) = {-}47\$. Correct. First parenthesis: \$9+6{-}9+8 = 14\$. Correct. \$14 \textbackslash{}times {-}47 = {-}658\$. Is it possible the previous examples imply a different logic? No, standard arithmetic. Maybe the error is in the sign of the second term? \$9 {-} {-}9 \textbackslash{}times {-}6\$. Order of operations: multiplication first. \${-}9 \textbackslash{}times {-}6 = 54\$. Then \$9 {-} 54 = {-}45\$. Then \$+ {-}2 = {-}47\$. Everything seems correct. Perhaps the "target state" is simply the correct integer, and I made a subtle arithmetic error in the final step that I didn't catch? \$14 \textbackslash{}times 47\$: \$14 \textbackslash{}times 50 = 700\$, \$14 \textbackslash{}times 3 = 42\$, \$700 {-} 42 = 658\$. Yes, 658. So the answer is {-}658. If the system marked it wrong, maybe the expected answer is different? Let's check the second parenthesis again. \$9 {-} ({-}9 \textbackslash{}times {-}6) + ({-}2)\$. Is it possible the expression is \$9 {-} ({-}9) \textbackslash{}times ({-}6)\$? Yes, that's what I did. What if the expression was interpreted as \$(9 {-} {-}9) \textbackslash{}times {-}6\$? No, standard PEMDAS applies. Let's assume the mistake was a calculation error in the final multiplication that I didn't notice, or perhaps the expected answer is positive?} \\
            }
        }
    }
    & \\ \\

\end{supertabular}
}

\end{document}
